\documentclass[a4paper,11pt]{article}
\usepackage{exam-style-341}
\examinehead{341 农业综合知识三 · 数据库技术与应用}{模拟试卷（二）}
\begin{document}
\examcover{341 农业综合知识三}{数据库技术与应用 · 模拟试卷（二）}{数据库技术与应用（50 分）}{50 分}{60 分钟}{本科目只考数据库，与程序设计、网络无关。}

\parttitle{数据库技术与应用（共 50 分）}

% ============================================================
\qtypename{一、单项选择题}{10 题，每题 2 分，共 20 分}
% ============================================================

\q 数据库系统的三级模式结构中，描述数据的物理存储结构的是（\quad）
\choicesline{A. 外模式}{B. 模式}{C. 内模式}{D. 视图}

\q 当数据库的物理存储结构改变时，通过修改（\quad）映像可以使模式保持不变，从而保证数据的物理独立性。
\choicesline{A. 外模式/模式}{B. 模式/内模式}{C. 外模式/内模式}{D. 外模式/存储模式}

\q 数据模型的三要素中，用于描述对数据允许执行的操作集合及其规则的是（\quad）
\choicesline{A. 数据结构}{B. 数据操作}{C. 完整性约束}{D. 数据联系}

\q 在关系 $R$ 中，属性（组）$X$ 不是 $R$ 的主码，但 $X$ 是另一个关系 $S$ 的主码，则 $X$ 在 $R$ 中称为（\quad）
\choicesline{A. 候选码}{B. 主属性}{C. 外码}{D. 超码}

\q 参照完整性规则的内容是（\quad）
\choices{主码的值不能为空值}{外码的值或者为空值，或者等于被参照关系中某个元组的主码值}{每个属性都必须是不可再分的}{关系中不允许出现完全相同的元组}

\q 与 3NF 相比，BCNF 进一步消除了（\quad）
\choices{非主属性对码的部分函数依赖}{非主属性对码的传递函数依赖}{主属性对码的部分函数依赖和传递函数依赖}{属性之间的多值依赖}

\q 设关系 $R(A,B,C)$ 和 $S(B,C,D)$，则 $R$ 与 $S$ 进行自然连接后，结果关系中包含的属性个数是（\quad）
\choicesline{A. 3}{B. 4}{C. 5}{D. 6}

\q 事务的原子性是指（\quad）
\choices{事务中包括的所有操作要么都做，要么都不做}{事务一旦提交，对数据库的改变就是永久的}{事务执行的结果必须使数据库从一个一致性状态变到另一个一致性状态}{一个事务的执行不能被其他事务干扰}

\q 下列关于视图的叙述中，\textbf{错误}的是（\quad）
\choices{视图是从一个或几个基本表导出的虚表}{视图可以嵌套定义，即可以在视图上再定义视图}{定义视图后，视图对应的数据会物理存储在数据库中}{视图能够对机密数据提供安全保护}

\q 关于两段锁协议，下列说法正确的是（\quad）
\choices{事务遵守两段锁协议是其并发调度可串行化的充分条件}{事务遵守两段锁协议就一定能避免死锁}{封锁粒度越大，系统的并发度越高}{共享锁允许其他事务修改被锁定的数据}

% ============================================================
\qtypename{二、填空题}{5 题，每题 2 分，共 10 分}
% ============================================================

\q 自然连接是一种特殊的等值连接，它要求两个关系中进行比较的分量必须是 \fillblank{} 的属性组，并在结果中去掉 \fillblank{}。

\q 在关系模式中，包含在任何一个候选码中的属性称为 \fillblank{}，不包含在任何候选码中的属性称为 \fillblank{}。

\q 关系数据库的三类完整性约束中，实体完整性和参照完整性由 DBMS 自动支持，而 \fillblank{} 完整性由用户根据具体应用环境自行定义。

\q 设关系模式 $R(A,B,C,D)$，函数依赖集 $F=\{A\to B,\ B\to C,\ C\to D\}$，则属性集 $AC$ 关于 $F$ 的闭包 $(AC)^{+}=$ \fillblank{}，$R$ 的候选码是 \fillblank{}。

\q SQL 是一种 \fillblank{}（填“过程性”或“非过程性”）语言，用户使用 SQL 时只需指出“做什么”，而不必指出“怎么做”。

% ============================================================
\qtypename{三、简答题}{2 题，每题 5 分，共 10 分}
% ============================================================

\q 什么是 E-R 模型（概念模型）？它在数据库设计中有什么作用？E-R 图中实体之间的联系有哪几种类型？
\anslines{3.2cm}

\q 数据库恢复的基本原理是什么？数据库系统可能发生的故障分为哪几类？简述各类故障的恢复方法。
\anslines{3.2cm}

% ============================================================
\qtypename{四、综合题}{2 题，每题 5 分，共 10 分}
% ============================================================

\q 某公司项目管理数据库中四个基本表的结构如下（加下划线的属性为主码）：

\begin{center}
\begin{tabular}{ll}
\toprule
关系名 & 属性（含义） \\
\midrule
Dept & \underline{Dno}, Dname, Dhead（部门号，部门名，负责人） \\
Emp & \underline{Eno}, Ename, Esex, Eage, Dno, Salary（职工号，姓名，性别，年龄，部门号，工资） \\
Project & \underline{Pno}, Pname, Budget, Dno（项目号，项目名，预算，承担部门号） \\
Work & \underline{Eno}, \underline{Pno}, Hours（职工号，项目号，工时） \\
\bottomrule
\end{tabular}
\end{center}

请完成下列各题：
\begin{enumerate}[label=（\arabic*）, leftmargin=2.4em, itemsep=1pt]
  \item（1 分）用 \code{CREATE TABLE} 语句定义基本表 \code{Work}，要求定义主码 \code{(Eno, Pno)}，并定义 \code{Eno}、\code{Pno} 分别到 \code{Emp}、\code{Project} 的外码，且当被参照表中的元组被删除时级联删除 \code{Work} 中相应的元组。
  \item（2 分）查询“研发部”（\code{Dname = '研发部'}）中工资高于本部门平均工资的职工的职工号、姓名和工资，结果按工资降序排列（要求使用相关子查询）。
  \item（2 分）查询参与了 2 个及以上项目、并且项目总工时超过 100 小时的职工的职工号、姓名、参与项目数和总工时，结果按总工时降序排列。
\end{enumerate}
\anslines{5.2cm}

\q 某高校排课系统中有一个关系模式

\begin{center}
$R$(学号 $Sno$, 课程号 $Cno$, 教师号 $Tno$, 教师名 $Tname$, 教室 $Room$)
\end{center}

\noindent 其语义为：一名学生选修一门课程只由一位教师任教；一位教师只有一个姓名；同一门课程由同一位教师任教时总在同一个教室。相应的函数依赖集为

\[
F=\{\,(\text{Sno},\text{Cno})\to\text{Tno},\quad
\text{Tno}\to\text{Tname},\quad
(\text{Tno},\text{Cno})\to\text{Room}\,\}
\]

请回答下列问题：
\begin{enumerate}[label=（\arabic*）, leftmargin=2.4em, itemsep=1pt]
  \item（1.5 分）写出 $R$ 的全部候选码，并指出其中的主属性和非主属性。
  \item（1.5 分）$R$ 最高属于第几范式？说明理由。
  \item（2 分）将 $R$ 分解为 3NF，要求分解既保持函数依赖又具有无损连接性，并用表格法（Chase 算法）验证该分解的无损连接性。
\end{enumerate}
\anslines{5.6cm}

\end{document}
